\uuid{RrlK}
\exo7id{3597}
\titre{exo7 3597}
\auteur{quercia}
\organisation{exo7}
\datecreate{2010-03-10}
\isIndication{false}
\isCorrection{true}
\chapitre{Réduction d'endomorphisme, polynôme annulateur}
\sousChapitre{Autre}
\cpge{mpsi-mp2i.calcul-matriciel-systemes.e, pcsi.determinants.c, ptsi.determinants.a, mpsi-mp2i.polynomes-fractions-rationnelles.a, pcsi.polynomes.a, ptsi.polynomes.a, bcpst1.matrices, mp-mpi.reduction-endomorphismes.b, pc.reduction.a, psi.reduction-endomorphismes.a, pt.reduction-endomorphismes.a, mp-mpi.reduction-endomorphismes.d, pc.reduction.c, psi.reduction-endomorphismes.c, pt.reduction-endomorphismes.c, mp-mpi.reduction-endomorphismes.g, pc.complements-espaces-vectoriels.d, psi.complements-espaces-vectoriels.d, bcpst2.valeurs-propres-vecteurs-propres.a, bcpst2.valeurs-propres-vecteurs-propres.b}
\module{Algèbre}
\niveau{L2}
\difficulte{}

\contenu{
\texte{
Soit $A \in \mathcal{M}_n(K)$ et $M = \begin{pmatrix} A &0 \cr A &A \cr\end{pmatrix} \in \mathcal{M}_{2N}(K)$.
}
\begin{enumerate}
    \item \question{Comparer les valeurs propres de $A$ et $M$.}
    \item \question{Soit $P \in  K[X]$ et $Q = XP'$.
    Montrer que $P(M) = \begin{pmatrix} P(A) & 0 \cr Q(A) & P(A) \cr \end{pmatrix}$.}
    \item \question{A quelle condition sur $A$, $M$ est-elle diagonalisable ?}
\reponse{
Par récurrence pour $P = X^k$, puis par linéarité.
$A=0$.
}
\end{enumerate}
}
